% sections/02_dataset_abstraction.tex  —  Dataset Abstraction
% \input{} into main.tex; do NOT wrap in \begin{document}.

\section{Dataset Abstraction}
\label{sec:dataset}

This section documents the structure, content, and integrity of the
available measurement data.  Precise characterisation of the dataset is a
prerequisite for every downstream decision—from representation design
(\Cref{sec:repr}) to evaluation protocol—because silent assumptions
about file format, frequency alignment, or class balance can propagate as
systematic errors through the entire pipeline.

% ─────────────────────────────────────────────────────────────────────
\subsection{File Format: Touchstone \texttt{.s1p}}
\label{sec:dataset:format}

The Touchstone format is the de facto standard for exchanging
frequency-domain network-parameter data between microwave instruments and
simulation tools~\cite{touchstone_spec, keysight_s2p_format}.  A
compliant Touchstone file begins with comment lines (prefixed by~\texttt{!})
followed by an \emph{option line} of the form

\begin{verbatim}
    # <freq_unit> <param> <format> <R> <ref_impedance>
\end{verbatim}

\noindent
which declares the frequency unit, parameter type ($S$, $Y$, $Z$, etc.),
data format (DB, MA, or RI), and reference impedance.  For a single-port
\texttt{.s1p} file the data section then contains one row per frequency
point with three columns: frequency, and two format-dependent values
(e.g.\ magnitude in~dB and phase in degrees for the DB format).

\subsubsection{Observed Deviation from the Standard}
\label{sec:dataset:format:deviation}

The \num{1603} files in this dataset \emph{lack} the standard
\texttt{\#\,\ldots} option header line.  Each file begins directly with
numeric data in a three-column, tab-separated layout with no comment
lines and no header row.  While the data content is consistent with the
Touchstone RI (Real-Imaginary) format, the absence of the option line means that:

\begin{enumerate}
    \item The frequency unit, data format, and reference impedance must
        be \emph{inferred} rather than parsed from metadata.
    \item Standard Touchstone parsers (e.g.\ \texttt{scikit-rf}) may
        reject or misinterpret these files without a preprocessing step
        that injects the missing header.
\end{enumerate}

\noindent
A robust ingestion pipeline should therefore prepend a synthetic option
line—e.g.\ \texttt{\# HZ S RI R 50}—before passing files to any
standards-compliant reader.

% ─────────────────────────────────────────────────────────────────────
\subsection{Column Interpretation}
\label{sec:dataset:columns}

Each data row contains three tab-separated fields:

\begin{center}
\begin{tabular}{cll}
    \toprule
    \textbf{Column} & \textbf{Quantity} & \textbf{Unit} \\
    \midrule
    1 & Frequency              & \si{\hertz} \\
    2 & $\text{Re}[S_{11}]$ (Real part) & linear / unitless \\
    3 & $\text{Im}[S_{11}]$ (Imag part) & linear / unitless \\
    \bottomrule
\end{tabular}
\end{center}

\noindent
All three quantities are stored as decimal floating-point numbers.
Frequency values include sub-hertz fractional digits (e.g.\
\texttt{8000999750.000}), reflecting the instrument's internal frequency
synthesis resolution.

% ─────────────────────────────────────────────────────────────────────
\subsection{Sample Trace Inspection}
\label{sec:dataset:sample}

\Cref{tab:sample_trace} reproduces the first five and last three rows of
the file \texttt{Trc0.s1p} from the \texttt{dry sensor} folder.  This
trace is representative of the general file structure across all
\num{1603} samples.

\begin{table}[htbp]
    \centering
    \caption{%
        First five and last three rows of \texttt{Trc0.s1p}
        (\texttt{dry} class) assuming RI format. Frequency is in~\si{\hertz}, 
        followed by Real and Imaginary linear components.%
    }
    \label{tab:sample_trace}
    \small
    \begin{tabular}{rrrr}
        \toprule
        \textbf{Row} &
        \textbf{Frequency (\si{\hertz})} &
        \textbf{Real ($\text{Re}[S_{11}]$)} &
        \textbf{Imaginary ($\text{Im}[S_{11}]$)} \\
        \midrule
        1    & 8\,000\,000\,000.000 & $0.065$ & $-0.013$ \\
        2    & 8\,000\,999\,750.000 & $0.066$ & $-0.012$ \\
        3    & 8\,001\,999\,500.000 & $0.063$ & $-0.010$ \\
        4    & 8\,002\,999\,250.000 & $0.065$ & $-0.006$ \\
        5    & 8\,003\,999\,000.000 & $0.062$ & $-0.006$ \\
        \midrule
        3999 & 11\,997\,000\,750.000 & $0.098$ & $-0.007$ \\
        4000 & 11\,998\,000\,500.000 & $0.100$ & $-0.003$ \\
        4001 & 11\,999\,000\,250.000 & $0.099$ & $-0.011$ \\
        \bottomrule
    \end{tabular}
\end{table}

Several observations follow from \Cref{tab:sample_trace}:

\begin{itemize}
    \item The frequency grid begins at exactly \SI{8.0}{\giga\hertz} and
        ends at approximately \SI{12.0}{\giga\hertz}
        (\SI{11999000250}{\hertz}), with a nominal step of
        $\Delta f \approx \SI{999750}{\hertz} \approx \SI{1}{\mega\hertz}$.
    \item The step size is not an exact integer number of hertz,
        consistent with a VNA synthesiser dividing a
        \SI{3.999}{\giga\hertz} span into exactly \num{4000} equal
        intervals.
    \item Real and Imaginary components span $[-1, 1]$ bounds, typical
        of a passive reflective structure.
    \item Because the representation is Cartesian, the curves are mathematically continuous;
        there are no ``phase wrap'' discontinuities requiring heuristic unwrapping logic.
\end{itemize}

% ─────────────────────────────────────────────────────────────────────
\subsection{Frequency Grid Validation}
\label{sec:dataset:freqgrid}

A uniform frequency grid is a critical assumption for all downstream
feature representations.  If any file deviates in the number of points or
in the frequency values themselves, sample-wise operations
(concatenation, differencing, normalisation) would silently produce
misaligned features.

The expected grid is defined by three parameters:
\begin{align}
    f_{\text{start}} &= \SI{8\,000\,000\,000}{\hertz},\notag\\
    f_{\text{stop}}  &= \SI{11\,999\,000\,250}{\hertz},\notag\\
    N &= 4001 \text{ points},
    \label{eq:freq_grid}
\end{align}
yielding a step of $\Delta f = (f_{\text{stop}} - f_{\text{start}}) /
(N-1) \approx \SI{999750.0625}{\hertz}$.  The ingestion pipeline must
verify, for every file, that (i)~the number of rows equals \num{4001},
and (ii)~the frequency column matches the reference grid to within a
tolerance of $\pm\SI{1}{\hertz}$.  Files failing either check should be
flagged and excluded rather than silently truncated or interpolated.

% ─────────────────────────────────────────────────────────────────────
\subsection{Class Distribution}
\label{sec:dataset:classes}

\Cref{tab:class_dist} summarises the number of traces per class and
quantifies the imbalance.

\begin{table}[htbp]
    \centering
    \caption{%
        Class distribution across the three surface-state categories.
        The imbalance ratio is computed relative to the minority classes.%
    }
    \label{tab:class_dist}
    \begin{tabular}{lrrr}
        \toprule
        \textbf{Class} & \textbf{Folder name} & \textbf{Traces} &
            \textbf{Ratio} \\
        \midrule
        \texttt{dry} & \texttt{dry sensor} & 601 & 1.20 \\
        \texttt{wet} & \texttt{wet sensor} & 501 & 1.00 \\
        \texttt{ice} & \texttt{ice on sensor} & 501 & 1.00 \\
        \midrule
        \textbf{Total} & & \textbf{1603} & \\
        \bottomrule
    \end{tabular}
\end{table}

The imbalance ratio of $1.20\!:\!1\!:\!1$ is mild by the standards of
applied machine learning.  The \texttt{dry} class has \num{100}
additional traces relative to \texttt{wet} and \texttt{ice}, which are
equal in size.  At this level of imbalance, standard stratified splitting
and stratified cross-validation are typically sufficient to prevent
biased evaluation.  Aggressive oversampling or synthetic minority
generation (e.g.\ SMOTE) is unlikely to be necessary but remains
available as a contingency.

% ─────────────────────────────────────────────────────────────────────
\subsection{Trace Numbering and Missing Acquisitions}
\label{sec:dataset:gaps}

File names follow the pattern \texttt{Trc\textlangle{}N\textrangle{}.s1p}
where $N$ is an integer index.  Inspection of the full file listing
reveals two systematic gaps in the numbering sequence:

\begin{itemize}
    \item \textbf{Trc601--Trc699} (\num{99} indices) are absent across
        all class folders.
    \item \textbf{Trc1201--Trc1449} (\num{249} indices) are absent
        across all class folders.
\end{itemize}

These gaps are consistent across the three class directories, suggesting
a measurement-campaign-level event (e.g.\ instrument recalibration,
paused acquisition, or deliberate exclusion of a batch) rather than
per-class data loss.  The gaps should be noted in any publication
describing the dataset but do \emph{not} affect the machine learning
pipeline, which operates on the files that are present and assigns
canonical indices independently of the original trace numbers.

% ─────────────────────────────────────────────────────────────────────
\subsection{Folder-to-Label Mapping}
\label{sec:dataset:labels}

Class labels are derived directly from the folder hierarchy.  Each
measurement file resides in one of three directories whose names serve as
ground-truth labels:

\begin{center}
\begin{tabular}{ll}
    \toprule
    \textbf{Folder name} & \textbf{Assigned label} \\
    \midrule
    \texttt{dry sensor}    & \texttt{dry} \\
    \texttt{wet sensor}    & \texttt{wet} \\
    \texttt{ice on sensor} & \texttt{ice} \\
    \bottomrule
\end{tabular}
\end{center}

\noindent
There is no auxiliary label file; the folder membership \emph{is} the
label.  This simple scheme avoids label-file synchronisation errors but
means that any accidental file movement would silently corrupt the
ground truth.  The ingestion pipeline should therefore checksum the full
file listing against a frozen manifest at the start of every experiment.

% ─────────────────────────────────────────────────────────────────────
\subsection{Deterministic Indexing Scheme}
\label{sec:dataset:indexing}

To ensure reproducibility, we propose a canonical sample~ID that is
independent of the original trace numbering.  Each sample receives an
identifier of the form

\begin{center}
    \texttt{\textlangle{}class\_prefix\textrangle{}\_\textlangle{}zero-padded index\textrangle{}}
\end{center}

\noindent
where:
\begin{itemize}
    \item \textbf{class\_prefix} is one of \texttt{dry}, \texttt{wet},
        or \texttt{ice}.
    \item \textbf{zero-padded index} is a four-digit integer assigned by
        sorting the files within each class folder in \emph{lexicographic
        order} by their original file name, then numbering from~\texttt{0000}.
\end{itemize}

\noindent
For example, the first file in the \texttt{dry sensor} folder
(\texttt{Trc0.s1p}) receives the canonical~ID \texttt{dry\_0000}; the
last file in the \texttt{ice on sensor} folder receives
\texttt{ice\_0500}.  This scheme satisfies three requirements:

\begin{enumerate}
    \item \textbf{Uniqueness.}  Every sample has a globally unique
        identifier.
    \item \textbf{Determinism.}  The same set of files always produces
        the same mapping, regardless of operating-system directory
        traversal order, provided lexicographic sorting is used.
    \item \textbf{Traceability.}  The original file name is recorded
        alongside the canonical~ID in the dataset manifest, preserving
        the link to the raw acquisition.
\end{enumerate}

% ─────────────────────────────────────────────────────────────────────
\subsection{Metadata Logging Requirements}
\label{sec:dataset:metadata}

Every experimental run must record the following metadata to enable
post-hoc auditing and reproduction:

\begin{enumerate}
    \item \textbf{Dataset manifest:}  A table mapping each canonical
        sample~ID to its original file path and an integrity hash
        (e.g.\ SHA-256 of the raw file bytes).
    \item \textbf{Frequency grid reference:}  The exact $f_{\text{start}}$,
        $f_{\text{stop}}$, and~$N$ used, plus the tolerance applied
        during grid validation.
    \item \textbf{Class distribution:}  The number of samples per class
        after any exclusion or filtering step.
    \item \textbf{Exclusion log:}  A list of any files excluded during
        ingestion, with the reason for exclusion (e.g.\ wrong number of
        rows, frequency grid mismatch, corrupt data).
    \item \textbf{Software versions:}  Python interpreter version and
        versions of all critical libraries (NumPy, scikit-learn,
        TensorFlow/PyTorch, XGBoost, etc.).
    \item \textbf{Random seeds:}  All seeds used for splitting,
        initialisation, and any stochastic algorithm component.
    \item \textbf{Preprocessing parameters:}  Normalisation statistics
        (mean, standard deviation, min, max) computed on the training
        fold, along with any augmentation or feature-engineering
        settings.
\end{enumerate}

\noindent
These requirements are not optional conveniences; in a small-sample
regime where every split decision materially affects reported
performance, full provenance is the only defence against irreproducible
results.
